一个配送时长预测模型上线八个月之后被下掉了。

复盘会开了三个小时，四个团队给出了四种解释。算法团队拿出监控：预测误差从上线到下线一直稳定，模型没有退化。数据团队指出训练集里全是成功送达的订单，被取消的那些从来没有进过样本。产品团队翻出最初的需求文档，上面写的是``提升用户对配送时间的满意度''，而实际交付的是一个把绝对误差最小化的回归模型——中间那次翻译谁也说不清是哪次会议上完成的。运营团队最后补了一句最关键的：这个预测值不只是显示给用户看，它同时被派单系统当作时间预算用；于是模型报得越乐观，骑手的时间就越紧，越紧就越容易跳过等待电梯这类耗时环节，实际送达时间因此真的变短了——直到骑手大量流失。

四种解释都对，而且它们不冲突。它们分别属于这条链路上的四个不同环节，症状在最后一个环节暴露，原因散落在前面三个。

读完这本书前十六个部分之后，最容易剩下的困惑是：每一章都懂了，可它们合起来在回答什么问题。这个单元给出的回答是走一遍上面那条链路——从有人提出一个诉求开始，到系统上线之后开始改变它所预测的那个世界为止，每一站说清三件事：这一步真正的难点是什么，判断依据在哪一章，以及这一步做错时故障会以什么形态出现在下游。

第三件事是这个单元的核心。现实中先看到的永远是症状，然后才需要反查是哪一环出的问题；而症状与病因几乎从不在同一站。

\begin{center}
\begin{tikzpicture}[x=1cm,y=1cm,font=\small,>=stealth]
  \draw[black!45,fill=blue!8] (1.2,5.83) rectangle (5.0,6.58);
  \node at (3.1,6.2) {一\ \ 目标与约束};
  \draw[black!45,fill=blue!8] (1.2,4.78) rectangle (5.0,5.53);
  \node at (3.1,5.15) {二\ \ 数据与表示};
  \draw[black!45,fill=blue!8] (1.2,3.73) rectangle (5.0,4.48);
  \node at (3.1,4.1) {三\ \ 模型族};
  \draw[black!45,fill=blue!8] (1.2,2.68) rectangle (5.0,3.43);
  \node at (3.1,3.05) {四\ \ 训练};
  \draw[black!45,fill=blue!8] (1.2,1.63) rectangle (5.0,2.38);
  \node at (3.1,2.0) {五\ \ 评估};
  \draw[black!45,fill=orange!12] (1.2,0.58) rectangle (5.0,1.33);
  \node at (3.1,0.95) {六\ \ 部署与反馈};
  \foreach \y in {5.83,4.78,3.73,2.68,1.63} \draw[->,black!55] (3.1,\y) -- (3.1,\y-0.3);
  \node[anchor=west,font=\footnotesize,black!55] at (5.6,6.85) {下游看到的典型症状};
  \node[anchor=west,font=\footnotesize,black!70] at (5.6,6.2) {指标全涨而体验变差};
  \node[anchor=west,font=\footnotesize,black!70] at (5.6,5.15) {离线与线上稳定差一截};
  \node[anchor=west,font=\footnotesize,black!70] at (5.6,4.1) {怎么调都卡在某个水平};
  \node[anchor=west,font=\footnotesize,black!70] at (5.6,3.05) {平台期、换种子结果大变};
  \node[anchor=west,font=\footnotesize,black!70] at (5.6,2.0) {反复不显著又反复上线};
  \node[anchor=west,font=\footnotesize,black!70] at (5.6,0.95) {先涨后跌、越重训越极端};
  \draw[->,red!70!black,very thick] (1.2,0.95) -- (0.45,0.95) -- (0.45,5.15) -- (1.15,5.15);
  \node[rotate=90,anchor=south,font=\footnotesize,red!70!black] at (0.2,3.0) {上线改写了训练数据};
\end{tikzpicture}

\emph{六站按时间顺序往下走，症状却几乎总是出现在病因的下游。红色那条回边是第六站独有的：它把一条流水线变成了一个会自己演化的系统，前五站的分析全都默认它不存在。}
\end{center}

\subsection{六站各自承诺了什么}

\hyperref[sec:17-Synthesis/01-System-Lifecycle/01]{``把诉求写成目标函数与约束''}处理最容易被跳过的一步。业务方说的话不是目标函数，把它变成目标函数的过程中做了一连串没被记录的选择：选哪个可测量的代理、约束写成硬约束还是罚项、多个目标怎样加权。这一步的核心判断是目标函数会被精确地钻空子，代理与真实目标的偏离在优化压力下会被系统性放大——这是结构性的，不是执行不力。

\hyperref[sec:17-Synthesis/01-System-Lifecycle/02]{``数据与表示里被悄悄用掉的假设''}拆开``数据是客观的''这个印象。采样机制决定了手上的分布究竟是什么，特征工程把一部分先验硬编码了进去，而标准化、缺失填充、特征选择这些看似中性的预处理本身携带信息——放在划分之前做，验证集的信息就已经进了训练。表示的选择同时是度量的选择，而高维下距离会集中、样本协方差的谱会失真，朴素的几何直觉在这里靠不住。

\hyperref[sec:17-Synthesis/01-System-Lifecycle/03]{``选模型族就是选一组关于世界形状的假设''}把选型问题从``哪个更强''改写成``哪组假设更接近这份数据''。线性假设效应可加，树假设边界轴对齐，核把假设写进核函数，概率图把假设写成哪些依赖不存在，卷积假设平移等变，注意力假设任意位置可以直接对话。假设的强度与样本量是一条权衡曲线：强偏置在数据少时是优势，数据多时变成天花板。

\hyperref[sec:17-Synthesis/01-System-Lifecycle/04]{``训练能保证什么不能保证什么''}处理最容易产生错误安全感的一步。优化保证的是在这个目标函数下找到了一个好点，它对目标函数写得对不对一无所知；凸问题保证局部即全局，非凸问题连这个也没有；而数值上收敛不等于数学上收敛，步长、条件数与浮点精度都能让一条正确的算法给出错误的答案。

\hyperref[sec:17-Synthesis/01-System-Lifecycle/05]{``评估是在为一句关于未见数据的断言背书''}回到全书重复最多的那条判断：置信度属于程序，不属于这次的参数。被评估的对象是整条流程而不是末端那个模型，多次比较与边看边停会把名义错误率悄悄推高，低功效下通过显著性筛子的效应量被系统性夸大。泛化界极其宽松，它的用处是指出哪些因素影响泛化，不是预测泛化误差的数值。

\hyperref[sec:17-Synthesis/01-System-Lifecycle/06]{``上线之后系统开始改变它所预测的世界''}是这个单元最重要的一篇，因为它补上了前五步共同缺失的维度。部署把预测问题变成了反馈系统：推荐影响用户看到什么，于是影响下一批训练数据；风控拒绝的申请永远不会产生标签；预测拥堵的导航把车导向另一条路，于是那条路堵了。这不是世界自己在漂移，是模型造成的漂移，此时必须换用因果的语言，并且要认真对待一个控制论的老问题——把自己的输出喂回输入的系统会不会震荡或发散。

\subsection{这个单元怎么读}

这六篇不是内容摘要，也不是目录的另一种排法。它们不复述任何一章的推导，只做一件事：在每个真实的决策点上，指出该用哪一条判断，并把那条判断本身用一两句话说清楚，让没读过那一章的人也能懂，读过的人被唤醒。想深入某一条，顺着链接回到对应的章去。

按顺序读一遍是最省事的路径，因为每一篇的结尾都交代了下一篇要补上什么。手上正卡着一个具体问题时，另一种用法是从图右侧那一列症状进入：先找到与手上现象最接近的那一条，读对应的一篇，然后往上游追。故障形态在每一篇末尾单独成节，就是为了支持这种逆向的用法。

需要提醒的是，这条链路呈现为六个先后的阶段，只是叙述上的方便。实际项目里它们互相纠缠——评估设计会反过来改变目标的写法，部署反馈会推翻当初的数据假设。把它写成一条线，是为了让每一站的责任边界清楚；真正做的时候，重要的是知道任何时刻手上这个问题属于哪一站，因为不同站点的问题需要完全不同的动作，而混淆它们是绝大多数无效努力的来源。
